\documentclass[11pt,a4paper]{article}
\usepackage[T1]{fontenc}\usepackage[utf8]{inputenc}\usepackage{lmodern}
\usepackage[a4paper,margin=2.2cm]{geometry}
\usepackage{amsmath,amssymb}\usepackage{enumitem}\usepackage{booktabs,longtable,array}
\usepackage[hidelinks]{hyperref}
\setlength{\parindent}{0pt}\setlength{\parskip}{0.4em}
\title{\textbf{OP-grad G2: are $Z_B^{(0)},Z_W^{(0)}$ frozen matching constants
or local functionals of $u_{\rm cos}$?}\\\large Recon v1 --- proposal-only;
target $\partial\ln Z^{(0)}/\partial\ln u_{\rm cos}=0$ at fixed bare data, or
an explicit limitation}
\author{Per reviewer directive after G1 acceptance}
\date{12 June 2026 --- rev.~1.1 (review safeguards applied)}
\begin{document}\maketitle

\section*{0. Scope and firewall}
G2 derives or classifies the $u$-dependence of the gauge-kinetic matching
coefficients. Outcome below: a \emph{conditional frozen-matching lemma} with
one condition holding at the level of the recorded gauge assignments, one
recorded-and-pinned condition, and one condition identical to a D2 hypothesis --- plus a
quantified counterfactual showing what a violation would cost. Proposal-only;
no preprint edits.

\section{Reformulation: ``frozen'' is an identity, not a memory effect}
Local EFT has no mechanism by which a Wilson coefficient ``remembers'' its
formation epoch: coefficients are either $u$-independent functions of bare
data, or they carry \emph{explicit} field dependence as operators (as the
gravitational kernel does through its $u(X)^2$ prefactor). The G2 question is
therefore exact and local: \emph{does the matching integral that produces
$Z_{B,W}^{(0)}$ contain a $u$-leg at fixed bare data?} ``Frozen matching''
holds iff it does not --- a property of the integral, not an initial
condition.

\section{Recorded-assignment input: the $u$-coupled heavy mode carries no
charge label}
In the recorded decomposition $\partial_A\Phi=u\,n_A$ with $n_An_A=1$, the
amplitude $u$ is by construction the \emph{invariant magnitude}: all
directional and phase content is carried by $n_A$ and by the compact phase
sector, while the emergent charges are holonomy/winding data of the compact
directions (eq:holonomy, $e\oint A=2\pi n$; charge lattice). A magnitude
fluctuation carries no winding and no orientation index. \textbf{In the
recorded decomposition, $\sigma=\delta u$ --- the only recorded $u$-coupled
heavy mode, the one integrated out in the stiffness kernel --- carries no
recorded holonomy or charge label; hence it is neutral with respect to the
recorded emergent gauge assignments, unless the microscopic completion adds
explicit portal operators} (condition (C-c) below).

\section{Consequence for the matching integrals}
One-loop gauge-kinetic coefficients are generated by vacuum polarisation of
\emph{charged} modes only: a neutral heavy field does not feed $F_{AB}F^{AB}$
at leading order. The $\sigma$-loop therefore feeds the orientation-gradient
operator $u^2(\partial n)^2$ --- kinematically, through the decomposition of
$|\partial\Phi|^2$, no charge required --- and the gravitational chain
$\kappa_n\to M_G^2$ (G1), \emph{but not} the gauge kinetic terms. The full
$u$-derivative of the matching coefficients at fixed bare data is then carried
by the charged spectrum alone:
\[
\frac{\partial\ln Z^{(0)}}{\partial\ln u_{\rm cos}}\Big|_{\rm bare}
=\sum_{f\,\in\,\text{charged}}\frac{k_f}{16\pi^2\,Z^{(0)}}\;
\gamma_{m_f^2},
\qquad k_f=O(1)\cdot q_f^2\,,
\]
through the threshold logarithms $\ln(\Lambda/m_f)$. The recorded charged
spectrum is mass-anchored to the electroweak chain, $m_f=y_f v_2$ with
$v_2=\phi_0e^{-\mathcal I_{\rm EW}}$, and the $\Phi$-sector evaluation point is
pinned at $(H/m_\varphi)^2\sim10^{-120}$ (D2), \emph{provided the
matter-sector matching point $u_0,\phi_0$ remains distinct from the
cosmological $u_{\rm cos}(z)$ excursion} (the D2 two-variable separation):
$\gamma_{m_f^2}\simeq0$.

\textbf{Conditional frozen-matching lemma.} Under
\begin{description}[leftmargin=2.0em,itemsep=2pt]
\item[(C-a)] $\sigma$ carries no recorded holonomy or charge label ---
\emph{neutral with respect to the recorded emergent gauge assignments} (\S2);
\item[(C-b)] all charged states have masses anchored to the pinned
$v_2$/$\phi_0$ chain, with no charged state of $u$-tracking mass ---
\emph{recorded as the inherited-SM mass chain; pinned by D2};
\item[(C-c)] no above-threshold direct $\sigma$--gauge portal operators ---
\emph{identical to D2 hypothesis H4, with the same $\gamma$-thresholds};
\end{description}
the matching integrals contain no $u$-leg and
\[
\boxed{\;\frac{\partial\ln Z_{B,W}^{(0)}}{\partial\ln u_{\rm cos}}
\Big|_{\rm bare}=0\;}
\]
--- frozen matching as a \emph{conditional} identity of the matching
integral, within the recorded charge assignments and the D2 separation
assumptions. The asymmetry against the
gravitational channel is structural, not tuned: \emph{gauge fields respond
only to charge; gravity responds to everything; the mode that drifts carries
no charge.}

\section{Quantified counterfactual: the cost of one violating state}
A single charged state with $u$-tracking mass ($\gamma_{m^2}=1$) contributes
\[
\gamma_Z\simeq\frac{k}{16\pi^2 Z^{(0)}}\approx\frac{1}{158\times4}
\approx1.6\times10^{-3},
\]
i.e.\ a factor $\sim30$--$50$ above the $3\times10^{-5}$ stress threshold
(under the stated stress-test assumptions). The verified anchors therefore
translate into a \emph{structural constraint on any completion}: every charged
state must satisfy $|\gamma_{m_f^2}|\lesssim2\times10^{-2}$ --- the charged
spectrum must be $u$-decoupled at the two-percent level. A new, falsifiable
selection rule exported to the completion programme.

\section{G2 verdict}
\begin{enumerate}[leftmargin=1.9em,itemsep=2pt,label=\textbf{V\arabic*.}]
\item The reviewer's target is met as a conditional lemma:
$\partial\ln Z_{B,W}^{(0)}/\partial\ln u_{\rm cos}=0$ at fixed bare data under
(C-a)--(C-c), with (C-a) holding at the level of the recorded gauge
assignments, (C-b) recorded
and D2-pinned, and (C-c) literally a D2 hypothesis. ``Frozen matching'' is a
conditional identity of the matching integral (no $u$-leg) within the stated
assumptions, not a memory mechanism.
\item The selectivity architecture acquires its structural origin:
charge-based mediation. The drifting amplitude carries no recorded charge
label; only charged
modes feed $F^2$; charged masses anchor to the pinned $\Phi$-sector. This is
the microscopic face of the D2 selective-dressing closure.
\item Explicit limitation (honesty clause): (C-b) is verified only against the
\emph{recorded} spectrum chain; a completion containing additional charged
states with $u$-tracking masses, or $\sigma$--gauge portals above the H4
thresholds, would re-open the channel at the quantified $10^{-3}$ level per
state. The lemma is conditional, not unconditional.
\item Per the review plan: with G1+G2 in hand, the next D1 candidate is the
upgrade of the OP-CONST block to a \emph{conditional selective-coupling
protection route supported by recorded closure + G1/G2} --- pending review of
this document. No preprint edits until then.
\end{enumerate}
\end{document}
